\section{ゴーストセル法による CCD 境界処理（代替実装）}
\label{app:ccd_ghost}
% ═══════════════════════════════════════════════════════════════════════════════

本付録は §\ref{sec:ccd_bc} の境界スキーム法に対する代替として，\textbf{ゴーストセル（ghost cell）法}
を説明する．本稿の標準実装は境界スキーム法であり（§\ref{sec:ccd_bc}），
第\ref{sec:pressure}章（§\ref{sec:ccd_neumann_bc}）の Neumann 境界条件の組み込みも境界スキーム法を前提としている．
独自実装としてゴーストセル法を選択する場合は，PPE の Neumann 境界処理（§\ref{sec:ccd_neumann_bc}）を
ゴーストセル法に合わせて再設計すること．

計算領域の外側に仮想セル（$i = -1$，$i = N$）を1層追加し，
境界条件に応じた値を設定することで，内点スキーム（3点ステンシル）を境界近傍でもそのまま適用できる．
これにより行列アセンブリのコードが統一される利点がある．

\begin{tcolorbox}[algbox, title={ゴーストセルの設定手順（左境界 $i=0$ の例）}]
\textbf{Step 1：ゴーストセル値の設定}

境界条件に応じて $f_{-1}$，$f'_{-1}$，$f''_{-1}$ を以下のように設定する：

\begin{center}
\renewcommand{\arraystretch}{1.4}
\begin{tabular}{lll}
\hline
境界条件 & ゴーストセル値 & 物理的意味 \\
\hline
Dirichlet（ノースリップ壁）: $u|_{\mathrm{wall}}=0$ &
  $u_{-1} = -u_1$ &
  鏡像反転（$u_{0} \approx 0$） \\
Dirichlet（強制流入）: $u|_{\mathrm{wall}}=U_w$ &
  $u_{-1} = 2U_w - u_1$ &
  $u_{0} = (u_{-1}+u_1)/2 = U_w$ \\
Neumann（流出・対称）: $\partial u/\partial x|_{\mathrm{wall}}=0$ &
  $u_{-1} = u_1$ &
  勾配ゼロ（偶関数的反射） \\
\hline
\end{tabular}
\end{center}

微分値のゴーストセル設定（ノースリップ壁の例）：
$f'$ は偶関数的に反射（$f'_{-1} = f'_1$），$f''$ は奇関数的（$f''_{-1} = -f''_1$）．

\textbf{Step 2：境界点での内点方程式の適用}

\begin{equation}
  \alpha_1 f'_{-1} + f'_0 + \alpha_1 f'_1
  = \frac{a_1}{h}(f_1 - f_{-1}) + b_1 h(f''_1 - f''_{-1})
  \label{eq:ccd_ghost_eq}
\end{equation}
ここで $f_{-1}$，$f'_{-1}$，$f''_{-1}$ はStep 1 で設定したゴーストセル値．
全格子点で同一の係数行列 $\mathbf{A}_L, \mathbf{A}_R, \mathbf{B}$ を使えるため，行列アセンブリが単純化される．

\textbf{注意：} 圧力の全 Neumann 境界問題（零固有モード）については第\ref{sec:pressure}章の対処法を参照すること．
\end{tcolorbox}

\begin{tcolorbox}[algbox, title={PPE Neumann 境界条件のゴーストセル法による再設計（$\partial p/\partial n = 0$ の場合）}]
境界スキーム法（§\ref{sec:ccd_bc}，§\ref{sec:ccd_neumann_bc}）では，
第1行を $p'_0 = 0$ に直接置き換えることで Neumann 条件を強制した．
ゴーストセル法を選択した場合，この条件は\textbf{ゴーストセル値の偶関数的反射}として実装する：

\medskip
\textbf{Step 1：圧力ゴーストセルの設定（左境界 $i=0$ の例）}
\[
  p_{-1} = p_1, \qquad p'_{-1} = p'_1, \qquad p''_{-1} = p''_1
\]
（偶関数的反射：$\partial p/\partial x|_{x=0} = 0$ を意味する）

\textbf{Step 2：境界点での内点方程式の適用}

ゴーストセル値を代入した Equation-I（式~\eqref{eq:ccd_ghost_eq}）：
\[
  \alpha_1 p'_{-1} + p'_0 + \alpha_1 p'_1
  = \frac{a_1}{h}(p_1 - p_{-1}) + b_1 h(p''_1 - p''_{-1})
\]
$p_{-1} = p_1$ および $p''_{-1} = p''_1$ を代入すると，右辺が消え：
\[
  (1 + 2\alpha_1) p'_0 = 0 \quad\Longrightarrow\quad p'_0 = 0
\]
Neumann 条件が自動的に成立する（境界スキームの明示的な修正は不要）．

\textbf{Step 3：CCD 演算子値の計算}

$p'_0 = 0$ が確定したうえで，変密度 PPE 演算子値
$(\mathcal{L}_{\mathrm{CCD}}^{\rho,x}\bm{p})_0 = p''_0/\rho_0$
（§\ref{sec:ccd_neumann_bc}，式の導出と同様；密度勾配項は $p'_0=0$ により消滅）を計算する．

\medskip
\textbf{零固有モードへの対処：} ゴーストセル法を使う場合も，全周 Neumann 境界での零モード問題（§\ref{sec:pinpoint}）は同様に発生する．式~\eqref{eq:pin_overwrite} の強制上書きによるピン条件が有効であり，ゴーストセル法特有の追加対処は不要である．
\end{tcolorbox}

% ═══════════════════════════════════════════════════════════════════════════════
